\subsection{Earlier bounds for matrix multiplication}\label{sec:history}

Strassen applied his seven-product identity recursively to blocks, giving
exponent $\log_2 7$~\cite{strassen}. Subsequent algorithms found savings
by working with several matrix products at once. Pan aggregated trilinear
forms~\cite{pan}, Bini showed how interpolation converts approximate
bilinear identities into exact algorithms~\cite{bini}, and Sch\"onhage's
asymptotic sum inequality turned a direct sum of matrix products into an
upper bound on the exponent of a single product~\cite{schoenhage}.
Coppersmith and Winograd studied recursive extraction of disjoint
products~\cite{cw82}. Strassen's laser method extracts them from tensor
powers~\cite{Strassen87}, and Coppersmith and Winograd designed an auxiliary
tensor for which progression-free sets accomplish the extraction~\cite{cw}.

The successive improvements depended on how many products could be kept
independent. Analyses of higher tensor powers~\cite{sto,ds,vw,legall},
refined extraction and asymmetric hashing~\cite{avw,duan,alphaomega,alman},
and numerical optimization~\cite{dupont} sharpened the bounds obtained
from the Coppersmith--Winograd tensor. Staggered extraction gave
$\omega_F<2.371055$ over every field~\cite{stagger}. A separate line of
work embedded matrix products in group algebras~\cite{cohn2003,cohn2005}
and in the algebras of coherent configurations~\cite{CohnUmansSupport}.

The upper bound $9/4$ comes from polynomial multiplication. In Strassen's
asymptotic spectrum~\cite{spectrum}, asymptotic rank is the maximum of
normalized functions that are additive on direct sums, multiplicative on
tensor products and monotone under restriction. Over $\mathbb C$, the
upper-bound proof tracks how these functions value polynomial multiplication
as the two input degrees grow. Degenerations give a lower bound on the
growth of this profile. Interpolation gives an upper bound that depends on
the value at matrix multiplication. Comparing the two bounds yields
$9/4$~\cite[Theorem~1.1]{nine}. The central degeneration separates blocks
that have independent variables on two legs and share the third. Partial
labels and semidirect sums give related formulations of this step
\cite{two258,li}. L\"u proved the separation step in arbitrary
characteristic~\cite{lu}. Amanov, van den Berg and Zuiddam obtained the
bound over every field~\cite[Theorem~1.1]{abvz}. See also the tensor-value
proof~\cite{tensorvalue} and Shpilka's explicit algorithms~\cite{shpilka}.
The rectangular version of the argument~\cite{uznanski,abvz} connects
these estimates to the dual exponent studied in~\cite{coppersmith,alphaomega}.

On the lower-bound side, Bl\"aser proved that the ordinary rank of $T_n$
is at least $\tfrac52 n^2-3n$ over every field~\cite{Blaeser}.
Over $\mathbb C$, Landsberg gave an ordinary-rank bound of
$3n^2-o(n^2)$~\cite{LandsbergRank}, and Landsberg and Micha{\l}ek
gave a border-rank bound of $2n^2-\lceil\log_2 n\rceil-1$
\cite{LandsbergMichalek}. A constant factor above $n^2$ disappears on
taking roots along the sequence $T_n^{\otimes j}\simeq T_{n^j}$.
Barriers for the laser method and bounds from tensor irreversibility
address algorithms that pass through specified intermediate
tensors~\cite{limits,irreversible}. For polynomial multiplication,
Amanov, van den Berg and Zuiddam similarly proved that $9/4$ is optimal
for a specified set of degenerations and combination
rules~\cite[Theorem~5.4]{abvz}. Theorem~\ref{thm:square} gives the
matching lower bound for unrestricted arithmetic circuits. Related
work on the spectrum constructs quantum functionals~\cite{CVZ} and
explicit points over arbitrary fields~\cite{almanlipratt}, and describes
the sublevel sets of asymptotic rank by polynomial
equations~\cite{SpectralSublevels}. Wigderson and Zuiddam survey this
framework~\cite{WZ}.
